\documentclass[11pt,a4paper]{article}
\usepackage[margin=1in]{geometry}
\usepackage{amsmath,amssymb,amsthm}
\usepackage{algorithm}
\usepackage{algpseudocode}
\usepackage{booktabs}
\usepackage{hyperref}

\theoremstyle{definition}
\newtheorem{definition}{Definition}[section]
\newtheorem{theorem}{Theorem}[section]
\newtheorem{lemma}[theorem]{Lemma}
\newtheorem{remark}{Remark}[section]

\title{\textbf{Rigorous Mathematical Specification, Spatial Subsampling, and Invariance Mechanics of Pooling Operators}}
\author{Team Scaibu \\ \small Email: \texttt{jaydeep.9664920749@gmail.com} \\ \small Architecture and Core Engine Team}
\date{October 2026}

\begin{document}
\maketitle

\begin{abstract}
This document presents the formal mathematical foundations, spatial subsampling formulations, and translation invariance characteristics of Pooling layers: Max Pooling, Average Pooling, and Global Average Pooling (GAP). Pooling layers reduce the spatial resolution of feature maps, control representation dimensionality, and inject local translation invariance into convolutional networks. We derive the analytical transformations, prove local shift invariance guarantees, analyze asymptotic operation complexities, trace a worked numerical example matching the reference implementation, and discuss modern strided convolution trade-offs.
\end{abstract}

\section{Notation and Mathematical Preliminaries}

Let $X \in \mathbb{R}^{C \times H_{\text{in}} \times W_{\text{in}}}$ denote an input feature activation tensor where $C$ is the channel count, $H_{\text{in}}$ is spatial height, and $W_{\text{in}}$ is spatial width.

\begin{itemize}
    \item $C \in \mathbb{N}_{\ge 1}$: Channel dimension (processed independently).
    \item $K_h, K_w \in \mathbb{N}_{\ge 1}$: Pooling kernel window dimensions.
    \item $s_h, s_w \in \mathbb{N}_{\ge 1}$: Spatial strides.
    \item $p_h, p_w \in \mathbb{N}_{\ge 0}$: Zero-padding margins.
    \item $H_{\text{out}}, W_{\text{out}} \in \mathbb{N}_{\ge 1}$: Output spatial grid resolution.
    \item $Y \in \mathbb{R}^{C \times H_{\text{out}} \times W_{\text{out}}}$: Output downsampled tensor.
\end{itemize}

\begin{table}[h]
\centering
\caption{Summary of Mathematical Symbols for Pooling Layers}
\begin{tabular}{lll}
\toprule
\textbf{Symbol} & \textbf{Type / Domain} & \textbf{Description} \\
\midrule
$X$ & $\mathbb{R}^{C \times H_{\text{in}} \times W_{\text{in}}}$ & Input feature map tensor \\
$K_h, K_w$ & $\mathbb{N}_{\ge 1}$ & Pooling receptive window size \\
$s_h, s_w$ & $\mathbb{N}_{\ge 1}$ & Subsampling stride factors \\
$Y_{\text{max}}$ & $\mathbb{R}^{C \times H_{\text{out}} \times W_{\text{out}}}$ & Max-pooled feature map \\
$Y_{\text{avg}}$ & $\mathbb{R}^{C \times H_{\text{out}} \times W_{\text{out}}}$ & Average-pooled feature map \\
$Y_{\text{GAP}}$ & $\mathbb{R}^{C \times 1 \times 1}$ & Global average-pooled channel representation \\
\bottomrule
\end{tabular}
\end{table}

\section{Problem Definition and Core Concepts}

\subsection{Intuitive Meaning}
Convolutional layers are translation-equivariant: shifting an object in an image shifts the corresponding activation map. For classification tasks, the network must become translation-invariant (recognizing an object regardless of precise pixel coordinates). Max pooling extracts the most salient feature within local neighborhoods, average pooling smooths local responses, and Global Average Pooling aggregates entire feature maps into a compact 1D vector per channel.

\subsection{Formal Mathematical Formulations}

\subsubsection{Max Pooling}
\begin{equation}
Y_{\text{max}}(c, h, w) = \max_{i=0}^{K_h-1} \max_{j=0}^{K_w-1} X_{\text{pad}}(c, h s_h + i, w s_w + j)
\end{equation}

\subsubsection{Average Pooling}
\begin{equation}
Y_{\text{avg}}(c, h, w) = \frac{1}{K_h K_w} \sum_{i=0}^{K_h-1} \sum_{j=0}^{K_w-1} X_{\text{pad}}(c, h s_h + i, w s_w + j)
\end{equation}

\subsubsection{Global Average Pooling (GAP)}
\begin{equation}
Y_{\text{GAP}}(c, 0, 0) = \frac{1}{H_{\text{in}} W_{\text{in}}} \sum_{h=0}^{H_{\text{in}}-1} \sum_{w=0}^{W_{\text{in}}-1} X(c, h, w)
\end{equation}

\section{Algorithm Specification}

\begin{algorithm}[H]
\caption{Generic Pooling Layer Forward Pass}
\label{alg:pooling}
\begin{algorithmic}[1]
\Require Input $X \in \mathbb{R}^{C \times H_{\text{in}} \times W_{\text{in}}}$, mode $M \in \{\text{max}, \text{avg}, \text{gap}\}$, window $(K_h, K_w)$, stride $(s_h, s_w)$, pad $(p_h, p_w)$.
\Ensure Output tensor $Y \in \mathbb{R}^{C \times H_{\text{out}} \times W_{\text{out}}}$, dimensions $H_{\text{out}}, W_{\text{out}}$.
\If{$M = \text{gap}$}
    \State Allocate $Y \in \mathbb{R}^{C \times 1 \times 1}$
    \For{$c \gets 0 \textbf{ to } C - 1$}
        \State $Y[c][0][0] \gets \frac{1}{H_{\text{in}} W_{\text{in}}} \sum_{h, w} X[c][h][w]$
    \EndFor
    \State \Return $\{Y, 1, 1\}$
\EndIf
\State $H_{\text{out}} \gets \lfloor (H_{\text{in}} + 2p_h - K_h) / s_h \rfloor + 1$, $W_{\text{out}} \gets \lfloor (W_{\text{in}} + 2p_w - K_w) / s_w \rfloor + 1$
\State Allocate $Y \in \mathbb{R}^{C \times H_{\text{out}} \times W_{\text{out}}}$
\For{$c \gets 0 \textbf{ to } C - 1$}
    \For{$h \gets 0 \textbf{ to } H_{\text{out}} - 1$}
        \State $h_0 \gets h \cdot s_h$
        \For{$w \gets 0 \textbf{ to } W_{\text{out}} - 1$}
            \State $w_0 \gets w \cdot s_w$
            \If{$M = \text{max}$}
                \State $Y[c][h][w] \gets \max_{i, j} X_{\text{pad}}[c][h_0 + i][w_0 + j]$
            \Else
                \State $Y[c][h][w] \gets \frac{1}{K_h K_w} \sum_{i, j} X_{\text{pad}}[c][h_0 + i][w_0 + j]$
            \EndIf
        \EndFor
    \EndFor
\EndFor
\State \Return $\{Y, H_{\text{out}}, W_{\text{out}}\}$
\end{algorithmic}
\end{algorithm}

\section{Theoretical Guarantees and Invariance}

\begin{theorem}[Local Translation Invariance of Max Pooling]
Let $X \in \mathbb{R}^{1 \times K \times K}$ be an isolated activation peak $X_{i^*, j^*} = M$ and $0$ elsewhere. For any internal spatial displacement $(\Delta i, \Delta j)$ such that $(i^* + \Delta i, j^* + \Delta j) \in [0, K-1]^2$:
\begin{equation}
\text{MaxPool}(T_\Delta X) = \text{MaxPool}(X) = M
\end{equation}
\end{theorem}

\begin{proof}
Because $\text{MaxPool}(X) = \max_{(i, j) \in \mathcal{W}} X_{ij}$, as long as the maximum activation coordinate $(i^*, j^*)$ remains inside the window $\mathcal{W}$, $\max_{(i, j)} (T_\Delta X)_{ij} = M$, guaranteeing exact invariance to local spatial perturbations.
\end{proof}

\subsection{Limitations}
\begin{itemize}
    \item \textbf{Information Loss}: Max pooling discards $75\%$ of spatial activations (for $2 \times 2$ pool), losing precise spatial coordinate localization required in segmentation.
\end{itemize}

\section{Complexity Analysis}

\subsection{Time Complexity}
\begin{itemize}
    \item \textbf{Comparison / Summation Operations}: Traverses $C \times H_{\text{out}} \times W_{\text{out}} \times K_h \times K_w$ elements.
    \item \textbf{Total Time Complexity}: $\Theta(C \cdot H_{\text{out}} \cdot W_{\text{out}} \cdot K_h \cdot K_w)$ operations.
\end{itemize}

\subsection{Space Complexity}
\begin{itemize}
    \item \textbf{Storage}: $\mathcal{O}(C \cdot H_{\text{out}} \cdot W_{\text{out}})$ memory for the downsampled output.
\end{itemize}

\section{Exactness and Error Analysis}

The reduction operators are exact floating-point comparisons and linear additions.

\section{Worked Numerical Example}

Let $C = 1, H_{\text{in}} = 2, W_{\text{in}} = 2, K = (2, 2), s = (2, 2), p = (0, 0)$.
\begin{equation}
X = [[[2.0, 8.0], [4.0, 6.0]]]
\end{equation}

\begin{enumerate}
    \item \textbf{Max Pooling}:
    \begin{equation}
    Y_{\text{max}} = [[[\max(2.0, 8.0, 4.0, 6.0)]]] = [[[8.0]]]
    \end{equation}
    \item \textbf{Average Pooling}:
    \begin{equation}
    Y_{\text{avg}} = \left[\left[\left[\frac{2.0 + 8.0 + 4.0 + 6.0}{4}\right]\right]\right] = \left[\left[\left[\frac{20.0}{4}\right]\right]\right] = [[[5.0]]]
    \end{equation}
    \item \textbf{Global Average Pooling}:
    \begin{equation}
    Y_{\text{GAP}} = [[[5.0]]]
    \end{equation}
\end{enumerate}

\section{Numerical Considerations and Edge Cases}

\begin{itemize}
    \item \textbf{Negative Value Padding in MaxPool}: Padding for MaxPool is initialized to $-\infty$ so that zero-padding boundaries do not artificially suppress negative activations.
\end{itemize}

\begin{thebibliography}{9}
\bibitem{boureau2010theoretical} Y.~L.~Boureau, J.~Ponce, and Y.~LeCun, ``A Theoretical Analysis of Feature Pooling in Visual Recognition,'' in \emph{International Conference on Machine Learning (ICML)}, pp.~111--118, 2010.
\bibitem{lin2013network} M.~Lin, Q.~Chen, and S.~Yan, ``Network In Network,'' in \emph{International Conference on Learning Representations (ICLR)}, 2014.
\end{thebibliography}

\end{document}
